\section{Data audit, preprocessing, and leakage controls}
\label{sec:audit}

\subsection{Q1a: Dataset audit}
\inputtable{q1_audit}
\inputtable{q1_anomalies}
\inputtable{q1_class_dist}
\inputtable{q1_duration_by_target}
\fig{q1_numeric_distributions}{Distributions of heavy-tailed / sentinel-valued numeric features (log count axis).}

% TODO(student): interpret
\studenttodo{Q1a: describe domain, target, feature types; discuss missing/\texttt{unknown} values
(esp.\ \texttt{poutcome}, \texttt{contact}), duplicates, anomalies (\texttt{pdays}=-1 sentinel, negative
\texttt{balance}, \texttt{campaign} extremes) and why \texttt{duration} is a problematic (leaking) predictor.}

\subsection{Q1b: Partitioning}
A stratified 70/15/15 train/validation/test split was created once with seed 42 before any fitting
(Table~\ref{tab:q1_split_sizes}); indices are stored in \texttt{data/splits/}. Partition sizes:
train \NTrain, validation \NVal, test \NTest. Stratified and plain random splitting were each repeated over 100
seeds on the full data and on a random 200-row sample (Table~\ref{tab:q1_split_comparison},
Figure~\ref{fig:q1_split_deviation}).
\inputtable{q1_split_sizes}
\inputtable{q1_split_comparison}
\fig{q1_split_deviation}{Absolute deviation of the positive-class proportion of each partition from the reference proportion, over 100 seeds.}

% TODO(student): interpret
\studenttodo{Q1b: justify the 70/15/15 + CV design; compare stratified vs.\ random deviations (full vs.\
small sample: zero/few positives in val/test); when can stratification itself fail (e.g.\ very rare classes,
temporal/grouped data, distribution shift, multi-label)?}

\subsection{Q1c: Leakage-safe preprocessing}
All transformations are steps of an sklearn/imblearn \texttt{Pipeline} (\texttt{src/preprocessing.py}) and are
therefore fit only on the training data or on the training folds of cross-validation. Numeric features: median
imputation, then standardisation (LogReg, kNN, GaussianNB) or quantile discretisation (BernoulliNB) or no
scaling (trees). Categorical features: missing values mapped to an explicit \texttt{"missing"} level, then
one-hot encoding with \texttt{handle\_unknown="ignore"} (unseen categories become an all-zero vector). The only
row-wise engineering (no fitted statistics) is the indicator \texttt{prev\_contacted} = [\texttt{pdays} $\neq -1$].
\inputtable{q1_transform_matrix}

% TODO(student): interpret
\studenttodo{Q1c: explain why each transformation matters (or not) for each model family (scale sensitivity of
distances/penalties, NB distributional assumptions, split invariance of trees); justify the constant
\texttt{"missing"} level instead of most-frequent imputation.}

\subsection{Q1d: Controlled leakage demonstration}
Model: Logistic Regression ($C=1$, \texttt{class\_weight=balanced}). Rows labelled \textbf{INVALID} are
deliberately wrong and are used only for this demonstration: INVALID\_1 includes \texttt{duration};
INVALID\_2 fits scaling and \texttt{SelectKBest}(k=20) on all development rows (train+validation) before CV;
INVALID\_3 adds the mean target of a high-cardinality group (job $\times$ month $\times$ day) computed on all
development rows including the row itself. The corresponding valid procedures fit the same steps inside the
pipeline (for target encoding, sklearn's cross-fitted \texttt{TargetEncoder}). The test set is not used.
\inputtable{q1_leakage_demo}
\fig{q1_leakage_demo}{Cross-validation and validation Macro-F1 of valid vs.\ deliberately INVALID procedures.}

% TODO(student): interpret
\studenttodo{Q1d: quantify the apparent improvement of each INVALID procedure over its valid counterpart and
explain why it is untrustworthy (information unavailable at prediction time / information from evaluation rows
entering the fit). Note which leakage was large and which was small here and why (e.g.\ global scaling with
31k rows changes little; target encoding of rare groups memorises labels).}

\subsection{Q1e: Protocol diagram and failure modes}
\fig{q1_protocol}{Experimental protocol.}

% TODO(student): interpret
\studenttodo{Q1e: explain at least three failure modes that invalidate results even if the code runs (e.g.\
preprocessing fit outside folds, selecting on test, leaking post-outcome features, duplicate/near-duplicate
rows across partitions, temporal ordering of campaigns, threshold tuned on evaluation data, unstable single
split / no variance estimate).}
